\chapter{The Electron Is Named}

\section{Naming makes it measurable}

This is the volume's climax, and it turns on a fact so plain you can check it by looking. The residue of the last
chapter is one object, read many ways. This chapter \emph{names} it --- and the surprise, the thing the whole book has
been building toward, is that naming is not a label stuck on after the measurement. Naming is what \emph{makes} the
object measurable. And the naming costs nothing: no axiom, no wish, no continuum. It is forced, for free, by a single
finite fact you can see with your own eyes.

Begin with the representation, because everything follows from its being finite. The machine sorts the formal terms of
the construction --- the tower of variations, of which there are endlessly many --- into a fixed, small number of
boxes. Not approximately; exactly. There are \emph{two} boxes. The constant term, the fixed point, goes in one; the
first variation and the second variation both go in the other. The type of a box is a finite type --- two elements, no
more --- and it is finite by construction, not by any assumption bolted on. To name a term is simply to compute which
box it falls in. The name \emph{is} the box, and there are only two of them. This is the whole apparatus: an endless
supply of formal terms, and a finite pair of boxes to sort them into.

Now the fact the climax turns on, and it is the oldest fact in combinatorics. Put more terms than there are boxes, and
two of them must land in the same box. There is no way around it: with one more term than boxes, some pair collides.
The machine proves exactly this. Take one more term than there are boxes and sort them; the construction shows --- not
by asserting it, but by checking every case and finding the last corner impossible --- that two distinct terms must
share a box. The proof is a handful of case splits and a decision procedure; the machine says of it, in its own
comment, that there is \emph{no continuum, no finite-cardinality machinery, and no choice} anywhere in it. And it
states the engine whole: an infinite tower of terms, sorted into finitely many boxes, cannot be laid out as a line ---
it must \emph{wrap}. The tower closes into a cycle, and the closing is forced by finiteness alone.

This is the moment the preface was pointing at. A choice among infinitely many things, the preface said, may need the
axiom of choice --- a wish granted, a selector posited where none can be built. But a choice among \emph{finitely} many
things needs no axiom at all. You can just look. And here, at the peak, that is exactly what happens: the name is not
chosen, not posited, not wished into being. It is \emph{forced}, by a finite count, in a proof with no choice in it ---
and the forcing is a thing you can watch happen. The whole book's discipline, the whole preface's promise, is cashed in
this one theorem: finiteness names the term, and finiteness needs no axiom.

And now why the forced name is a \emph{measurement}. To say that two terms share a box is to say something precise and
decidable: that the representation \emph{cannot tell them apart}. Box-sameness is indistinguishability --- the exact
word the first chapter began with --- and it is decided natively, computed and not assumed; the machine flags the
premises that establish it as decided facts, never axioms and never hypotheses. So a named box is finite, it is
decidable, and it is distinguishable from the other box --- which is to say it is precisely a \emph{reading}, in the
sense the whole book has meant by the word since its first page. The residue that separates the first variation from
the second falls below the representation's resolution: the two share a box, and finiteness names them the same. That
named box is the measurable object. Naming did not describe a measurement taken elsewhere; naming \emph{was} the
measurement --- the term resolved into a finite, decidable, distinguishable box, which is all a reading has ever been.

In the two verbs this is the climax, the pair's deepest appearance in the book. The first chapter \emph{tanged}: it
told each term apart from the others, selecting by the characteristic that distinguished it. But telling apart has a
limit, and finiteness is where it is reached. With more terms than boxes, the terms \emph{cannot} all stay apart ---
there are not enough boxes to keep them separate, and so they are forced to \emph{funge}, to collide into shared boxes
whether or not you wanted them to. The name is that compelled funge: two terms bagged as one, not because you chose to
gather them but because the finite representation could not tell them apart. Identity, at the peak of the book, is the
funge finiteness forces --- indistinguishable to the representation, and named the same for exactly that reason.

So this chapter opens on the deepest thing the volume has to say, and it is a small thing said exactly. A finite
representation, faced with infinitely many terms, must name two of them the same --- and the name it forces is the
reading. This is the apex the preface aimed at and the ending will rest on. When the book finally hands back its
bracket --- the interval floored at $\alpha$ --- that bracket is measurable for the very same reason the name is forced
here: because the machine's representation is finite, and a finite representation you can just look at, no axiom
required. What remains for this chapter is only to see which box the forced name lands in, and what the machine has, all
along, been calling it.

\section{The stable reading}

The last section forced the naming: the pigeonhole put the recovered term into a box, and no axiom was needed to do it.
This section reads that box. And the point of the reading is that it is \emph{stable} --- not because the machine
insists on it, but because it is \emph{forced}, and a forced reading is one you could not have made otherwise. In a
world of two boxes, the box the term lands in is settled the moment you know it is not the other one. There is nowhere
else for it to go.

Start with which box the forced name lands in. One box holds the constant term --- the fixed point, the value the
construction pins down first. The recovered term is not that. It is the surviving leading term of the variation, the
part that did not reduce to the constant, and the construction proves it directly: the box of the recovered term is
\emph{not} the value's box. So the named object is, precisely, the term that is \emph{not} the value --- the second
box, the one the value does not occupy. The whole content of the naming, so far, is negative: the machine has found the
term that is left over when the value is taken out.

Now why that reading is stable, and it is the heart of the section. With only two boxes, ``not the value's box'' has
exactly one meaning: the other box. The construction proves this too, and proves it by decision --- for any box, if it
is not the value's, then it is the electron's; there is, in the machine's own words, nowhere else to sit. The reading
is not chosen from among options. Finiteness leaves no options: once the term is known not to be the value, its box is
forced, and the proof that it is forced runs by the same choice-free decision the last section turned on. This is what
stability \emph{means} here. A reading you could have made differently is a choice, and a choice can be revised or
disputed. A reading you could not have made otherwise is forced, and a forced reading is stable --- there is no other
reading available to drift to. The name is stable because the finite count permits exactly one.

And here the book does something it has been careful about all along, and does it with the code's explicit permission.
It reads the forced box as a \emph{charge} --- as minus one, the electron's. But it marks that reading exactly, and the
mark is not the gauge's caution laid on from outside; it is written into the source itself. The comment at the forced
box says it plainly: the sign is the book's reading of this box; the machine's claim is the box's identity. Two claims,
and the code separates them for us. The \emph{proved} claim --- the one the machine checks --- is the box: that the
recovered term is the non-value box, forced there by finiteness. The \emph{read} claim --- the book's, not the
machine's --- is that this box's sign is minus one, the electron's charge. The construction proves the box; the book
reads the sign. And the reading has a source the volume defers: the sign comes from a doubled quarter-turn on the
imaginary axis a later chapter installs --- the reversal a rotation delivers when it is taken twice. That is the next
volume's physics and the next chapter's rotation; here the book stakes only what the code proves --- the forced integer
box --- and names its sign as a reading it will earn later.

One caution the book owes the reader at exactly this point, because it is the discipline the whole volume rests on. The
minus one named here is a \emph{discrete integer} --- the electron's charge, a whole number the pigeonhole forces and
the book reads off the forced box. It is emphatically \emph{not} the coupling. The fine-structure number the instrument
is built to bracket is a different thing entirely, and it stays blind to the very end of the book; naming the
electron's integer charge here costs the coupling's magnitude nothing. The volume stakes the integer it can prove and
leaves the coupling exactly where it has always been --- unnamed, un-numbered, waiting for its bracket.

In the two verbs the section closes the climax the last one opened. The naming was the compelled funge: finiteness
bagged the term into a box it could not avoid. This section reads that bagged box and finds the reading stable
\emph{because} the funge was compelled --- there was nowhere else to bag the term, so there is nowhere else for the
reading to go. Identity, here, is not a selection among boxes; it is the one box finiteness permits, read. And that is
the last thing the chapter needs before it can close: the electron is named, its box forced, its reading stable, and
its sign staked as exactly what it is --- a reading the book will earn, laid honestly beside a box the machine has
already proved.

\section{Named by counting, not fiat}

Two sections have named the electron and read the name stable. This closing section says the thing that makes the
naming \emph{honest}, and it is a claim about what was \emph{not} done. Nothing was posited at the naming. No axiom was
invoked, no premise added, no number decreed. The name is a composition of forces the construction had already proved
--- the pigeonhole that forces a collision, and the fact that the collided term is not the value --- and their
composition is the name, with nothing new laid on top. The electron is named \emph{by counting}, not by fiat.

Look at how the naming is actually proved, because the proof is a composition and composes nothing new. The theorem
that the recovered term is the electron's box is not established by a fresh argument; it is the earlier two theorems,
joined. The pigeonhole forced the collision; the not-value fact placed the collided term in the electron's box; and the
naming is exactly those two, applied in turn. The construction's own comment is exact about it: there is \emph{no new
axiom and no new premise at the identity} --- the name is checked off the steps already taken, and adds not one thing
to them. This is naming as the cheapest decision the machine can make, written down. It asks the pigeonhole --- the
plainest question it has, a question about a finite count --- and inscribes the answer as the name. Nothing is smuggled
in. The naming is free of fiat in exactly the way the very first difference was free of allocation: it costs nothing
beyond what was already there.

And now the structural reason a finite machine can name an infinite tower at all, which is the most beautiful fact in
the chapter. The formal terms of the construction go on forever --- the tower of variations has no last floor. But the
boxes do not: there are two. An endless sequence of terms sorted into two boxes cannot keep them all apart; it cannot
stay a straight line, one term to one box forever, because it runs out of boxes while the terms keep coming. So the
tower must \emph{fold}. Two distinct terms, somewhere up the tower, are forced into the same box --- and that forced
coincidence is not a failure of the sorting. It \emph{is} the name. The tower that cannot be laid out as a line closes,
instead, into a loop --- the same loop the early chapters first rehearsed, when a round-trip came back to where it
started. Here that loop is the naming itself: the tower, unable to go on forever as a line, comes back on itself, and
the coming-back is the identity the machine reads as the name.

There is a last gift in the closing, and it is the chapter's quiet payoff. Because the naming folds the tower into
finitely many boxes, it does something to the hardest term in it. The exact leading term of the construction --- the
full, nonlinear one --- and the linear term of the reduced problem land, the machine proves, in the \emph{same box}.
The proof is immediate; they are simply, by computation, the same. Read on the finite representation, the intractable
exact term and the tractable approximate one are indistinguishable: finiteness did not merely close the gap between
them, it flattened the hard problem onto the easy one. That flattening --- the reading that finiteness
\emph{linearizes} the exact problem --- is the book's, a gloss laid over the shared box, and the deeper machinery of it
belongs to the physical volume. What the code proves is smaller and exact: the two terms share a box. But the lesson is
the chapter's whole point restated. Naming by counting does not merely label the object; it turns a problem you could
not have solved exactly into a finite, decided one you can.

In the two verbs the climax ends exactly where the book began. The first chapter \emph{tanged}: it told terms apart.
But telling infinitely many terms apart requires infinitely many boxes, and there are two. So the tower cannot keep its
terms tanged; it \emph{must funge}, and the wrap is that forced funge closing on itself. Identity, at the end of the
climax, is counted and not decreed --- the bag finiteness leaves when it runs out of boxes, the loop the tower folds
into when it cannot go on as a line. The two verbs that opened the book close it: a tange the count funges, shown now
to be the only honest way a finite machine can name anything at all.

So the electron's chapter ends. The name was forced by a finite count, read stable because forced, staked as an integer
the machine proves and a sign the book reads, and closed, here, as a tower that could not be a line folding into the
loop that is its own name. Nothing was decreed; everything was counted. And that is the whole of what the volume has
claimed the electron to be --- not a thing posited, but a coincidence a finite representation could not avoid, named by
looking at it. What remains for the book is to rotate this named object onto a second axis, and to ask it, at last, the
one question it has been built to ask.
